% ===========================================================================
%  Kapitel 4 - Versuchsdesign und Benchmark-Methodik
%
%  Hinweise zur Uebernahme:
%   * Listings sind als \begin{program} ... \end{program} gesetzt, analog zu
%     den Programmen 3.1-3.4. Falls dein Template eine andere Umgebung
%     verwendet, ist nur diese zu ersetzen.
%   * Abbildung 4.1 verwendet TikZ. Dafuer \usepackage{tikz} in der Praeambel.
%   * Zitierschluessel sind Platzhalter und muessen an deine .bib angepasst
%     werden.
%   * Offene Punkte sind mit "% OFFEN:" markiert und am Dateiende gelistet.
% ===========================================================================

\chapter{Versuchsdesign und Benchmark-Methodik}

Dieses Kapitel beschreibt den Aufbau der Messungen. Es legt fest, welche
DSP-Bausteine untersucht werden, mit welchen Werkzeugen gemessen wird, auf
welcher Hardware die Messungen laufen, welche Compilerkonfigurationen
verwendet werden und welche Testdaten zum Einsatz kommen. Ziel des Aufbaus
ist ein reproduzierbarer Vergleich zwischen Faust-generiertem und manuell
implementiertem C++-Code. Die konkrete Umsetzung der Implementierungen wird
anschließend in Kapitel~5 beschrieben.


\section{Definition der untersuchten DSP-Bausteine}
\label{sec:bausteine}

\subsection{Auswahlkriterien}

Abschnitt~1.2 nennt zwei Anforderungen an die untersuchten Bausteine. Sie
sollen unterschiedliche Datenabhängigkeiten aufweisen und unterschiedliche
Skalierungseigenschaften besitzen. Beide Anforderungen bestimmen die Auswahl.

Die Datenabhängigkeit beschreibt, ob ein Ausgangssample von einem zuvor
berechneten Ausgangssample abhängt. Diese Eigenschaft entscheidet darüber, ob
der Compiler eine Schleife automatisch vektorisieren kann. Der Grund ist
einfach: Vektorisierung bedeutet, dieselbe Rechenoperation gleichzeitig auf
mehrere aufeinanderfolgende Samples anzuwenden. Das setzt voraus, dass diese
Samples unabhängig voneinander berechnet werden können. Hängt $y[n]$ von
$y[n-1]$ ab, muss zuerst $y[n-1]$ vollständig vorliegen. Die Werte müssen
dann nacheinander berechnet werden, und eine gleichzeitige Verarbeitung ist
nicht möglich. Abschnitt~2.4.1 beschreibt diesen Zusammenhang; die
Abschnitte~2.2.2 und~2.2.3 wenden ihn auf rekursive und nichtrekursive
Filter an. Der Einfluss der automatischen Vektorisierung ist zugleich eine
der Teilfragen aus Abschnitt~1.2.

Die Skalierungseigenschaft beschreibt, wie der Rechenaufwand mit einem
Parameter wächst. Untersucht werden drei solcher Parameter: die Blockgröße,
die Filterlänge und die Transformationslänge.

Auf dieser Grundlage werden fünf Bausteine untersucht.
Tabelle~\ref{tab:bausteine} fasst sie zusammen. Sie decken den Bereich von
vollständiger Datenparallelität bis zu vollständiger Rückkopplung ab.

\begin{table}[htb]
\centering
\caption{Untersuchte DSP-Bausteine und ihre Eigenschaften}
\label{tab:bausteine}
\begin{tabular}{llll}
\hline
Baustein & Datenabhängigkeit & Zustand & Variierter Parameter \\
\hline
Verstärkungsstufe & keine & keiner & Blockgröße 64--1024 \\
Akkumulator & Rückkopplung (1 Sample) & 1 Wert & Blockgröße 64--1024 \\
Biquad-Filter & Rückkopplung (2 Samples) & 4 Werte & Blockgröße 64--1024 \\
FIR-Filter & nur vorwärts & $L$ Werte & Filterlänge $L \in \{32, 64, 256\}$ \\
FFT & strukturell ausgerollt & $N$ Werte & Länge $N \in \{64, 128, 256\}$ \\
\hline
\end{tabular}
\end{table}

Die ersten vier Bausteine werden mit den Blockgrößen 64, 128, 256, 512 und
1024 Samples gemessen. Diese Größen entsprechen üblichen Puffergrößen in
Audioanwendungen. Die FFT bildet eine Ausnahme. Bei ihr ist die Blockgröße an
die Transformationslänge gebunden, weshalb nur die drei Längen 64, 128 und
256 gemessen werden. Der Grund für diese Begrenzung wird in
Abschnitt~\ref{subsec:fft} erläutert.

Der Akkumulator ist gegenüber der Aufzählung in Abschnitt~1.2 hinzugekommen.
Er wurde ergänzt, weil der Biquad-Filter Rückkopplung und
Koeffizientenrechnung miteinander vermischt. Der Akkumulator trennt beides
und misst die Rückkopplung allein.
% OFFEN: Den Akkumulator zusaetzlich in der Aufzaehlung in Abschnitt 1.2 und
% in Abschnitt 2.2 der Grundlagen ergaenzen, sonst taucht er hier erstmals auf.


\subsection{Verstärkungsstufe}

Die Verstärkungsstufe ist der einfachste Fall. Sie benötigt pro Sample genau
eine Multiplikation. Es gibt keinen Zustand und keine Abhängigkeit zwischen
den Samples. Jedes Ausgangssample kann unabhängig berechnet werden. Nach
Abschnitt~2.2.1 besitzt der Baustein damit das höchste
Vektorisierungspotenzial. Er dient als oberer Referenzfall.

\begin{program}
\caption{Faust-Beschreibung der Verstärkungsstufe}
\label{prog:gain}
\begin{lstlisting}
process = *(hslider("gain", 0.5, 0, 2, 0.01));
\end{lstlisting}
\end{program}

Die C++-Referenz multipliziert jedes Sample in einer Schleife mit einem
Faktor desselben Werts. Beide Implementierungen berechnen damit dasselbe
Ergebnis.

Die Wahl des Faktors ist nicht beliebig. Ein Faktor von 1 wäre das neutrale
Element der Multiplikation. Der Compiler erkennt das und entfernt die
Rechnung vollständig; die Schleife wird dann zu einer reinen Kopie. Eine
Prüfung des erzeugten Assemblercodes bestätigt dies: Bei einem Faktor von 1
enthält der Verarbeitungsteil ab der Optimierungsstufe \texttt{-O1} keinen
einzigen Multiplikationsbefehl mehr, und zwar bei beiden Implementierungen
gleichermaßen. Der Baustein würde in diesem Fall nicht die Verstärkung
messen, sondern den Speicherdurchsatz. Deshalb wird durchgehend ein Faktor
von 0{,}5 verwendet. Die Multiplikation bleibt dadurch im erzeugten Code
erhalten.
% OFFEN: Diese Beschreibung setzt voraus, dass der Faktor in dsp/gain.dsp,
% in src/native/gain_native.h UND in evaluation/reference_numpy.py auf 0.5
% geaendert und anschliessend neu gemessen wurde. Ohne die Aenderung in der
% NumPy-Referenz schlaegt der Genauigkeitstest fehl.


\subsection{Akkumulator}

Der Akkumulator bildet die Summe aller bisherigen Eingangswerte. Für ein
Eingangssignal $x[n]$ gilt

\begin{equation}
y[n] = x[n] + y[n-1].
\end{equation}

Der Baustein besitzt eine Rückkopplung über genau ein Sample. Er enthält
keine Filterkoeffizienten und keine weitere Logik. Damit isoliert er die
reinen Kosten einer Datenabhängigkeit. Er dient als unterer Referenzfall für
rekursive Strukturen und als Vergleichspunkt für den Biquad-Filter. In Faust
entspricht ihm der rekursive Kompositionsoperator aus Abschnitt~3.3.

\begin{program}
\caption{Faust-Beschreibung des Akkumulators}
\label{prog:accum}
\begin{lstlisting}
process = + ~ _;
\end{lstlisting}
\end{program}


\subsection{Biquad-Filter}

Der Biquad-Filter setzt die Differenzengleichung~(2.11) aus
Abschnitt~2.2.2 um. Verwendet werden die Koeffizienten $b_0 = 1$,
$b_1 = b_2 = 0$, $a_1 = -1{,}5$ und $a_2 = 0{,}6$. Der Vorwärtszweig ist
damit auf einen einzelnen Faktor reduziert. Der Baustein misst deshalb
ausschließlich die Kosten der Rückkopplung über zwei Samples. Die Pole des
Filters liegen bei einem Betrag von etwa $0{,}77$ und damit innerhalb des
Einheitskreises. Der Filter ist somit stabil, und der Zustand bleibt über
beliebig viele Blöcke endlich.

\begin{program}
\caption{Faust-Beschreibung des Biquad-Filters}
\label{prog:biquad}
\begin{lstlisting}
import("stdfaust.lib");
process = fi.tf2(1.0, 0.0, 0.0, -1.5, 0.6);
\end{lstlisting}
\end{program}

Die C++-Referenz implementiert dieselbe Gleichung in Direktform~I. Sie
speichert die beiden vergangenen Eingangswerte und die beiden vergangenen
Ausgangswerte als Zustand.


\subsection{FIR-Filter}

Der FIR-Filter besitzt keine Rückkopplung. Jedes Ausgangssample hängt nur von
aktuellen und vergangenen Eingangswerten ab, wie in Gleichung~(2.12)
beschrieben. Der Aufwand pro Ausgangssample wächst linear mit der
Filterlänge. Die Filterlänge ist deshalb der interessanteste
Skalierungsparameter der Arbeit.

Untersucht werden drei Filterlängen: 32, 64 und 256 Koeffizienten. Alle drei
Längen sind Zweierpotenzen. Diese Einschränkung folgt aus der eigenen
C++-Referenz. Sie speichert die vergangenen Eingangswerte in einem
Ringpuffer. Der Zugriff auf einen zurückliegenden Wert erfordert eine
Umrechnung des Index auf den gültigen Bereich. Bei einer Zweierpotenz genügt
dafür eine bitweise Und-Verknüpfung mit $L-1$. Andernfalls wäre eine
Divisionsrestoperation nötig, die deutlich teurer ist. Die Implementierung
prüft diese Bedingung zur Übersetzungszeit. Abschnitt~5.1 zeigt den
zugehörigen Quelltext.

Als Koeffizienten dient eine normierte Rampe:

\begin{equation}
b_k = \frac{2\,(k+1)}{L\,(L+1)}, \qquad k = 0, \dots, L-1 .
\end{equation}

Die Summe aller Koeffizienten beträgt genau 1. Der Filter besitzt damit bei
Gleichanteil eine Verstärkung von 1. Die Koeffizienten sind für alle Längen
nach derselben Formel gebildet. Dadurch bleiben die Ergebnisse zwischen den
drei Filterlängen vergleichbar.

\begin{program}
\caption{Faust-Beschreibung des FIR-Filters (hier mit $L = 64$)}
\label{prog:fir}
\begin{lstlisting}
import("stdfaust.lib");
N = 64;
process = fi.fir(par(i, N, 2.0*(i+1)/(N*(N+1))));
\end{lstlisting}
\end{program}


\subsection{Fast-Fourier-Transformation}
\label{subsec:fft}

Die FFT unterscheidet sich strukturell von den übrigen Bausteinen. Ihre
Faust-Beschreibung besteht aus zwei Teilen, die
Abbildung~\ref{fig:fftfluss} darstellt.

\begin{program}
\caption{Faust-Beschreibung der FFT (hier mit $N = 64$)}
\label{prog:fft}
\begin{lstlisting}
import("stdfaust.lib");
N = 64;
process = an.rtocv(N) : an.fft(N);
\end{lstlisting}
\end{program}

Der erste Teil, \texttt{an.rtocv(N)}, wandelt ein reelles Signal in einen
komplexen Vektor um. Dazu wird das Eingangssignal über eine
Verzögerungsleitung mit $N$ Abgriffen geführt. Die $N$ zuletzt eingetroffenen
Samples liegen dadurch gleichzeitig vor. Jeder Abgriff wird zu einem
komplexen Wert mit dem Sample als Realteil und dem Imaginärteil null. Die
Reihenfolge beginnt beim neuesten Sample.

Der zweite Teil, \texttt{an.fft(N)}, transformiert diese $N$ komplexen Werte
in $N$ Spektralwerte. Der Baustein besitzt damit einen Eingang und $2N$
Ausgänge. Die Ausgänge enthalten abwechselnd Real- und Imaginärteil eines
Koeffizienten.

\begin{figure}[htb]
\centering
\begin{tikzpicture}[
  node distance=6mm,
  box/.style={draw, rounded corners=2pt, minimum height=8mm, align=center,
              font=\footnotesize, inner sep=2mm},
  ar/.style={-latex}
]
\node[box, minimum width=20mm] (tap) {Verzögerungs-\\leitung, $N$ Abgriffe};
\node[box, minimum width=22mm, right=10mm of tap] (rtocv)
     {\texttt{an.rtocv(N)}\\$N$ komplexe Werte};
\node[box, minimum width=22mm, right=10mm of rtocv] (fft)
     {\texttt{an.fft(N)}\\$\log_2 N$ Stufen};
\node[right=10mm of fft, font=\footnotesize, align=center] (out)
     {$2N$ Ausgänge\\$\mathrm{Re}(X_0), \mathrm{Im}(X_0), \dots$};
\node[left=8mm of tap, font=\footnotesize] (in) {$x[n]$};

\draw[ar] (in) -- (tap);
\draw[ar] (tap) -- (rtocv);
\draw[ar] (rtocv) -- (fft);
\draw[ar] (fft) -- (out);

\node[below=7mm of rtocv, font=\footnotesize, align=center, text width=75mm]
     {Der gesamte Ablauf wird für \emph{jedes} Sample erneut ausgeführt.\\
      Ein Block von $N$ Samples erfordert daher $N$ Transformationen.};
\end{tikzpicture}
\caption{Signalfluss der Faust-FFT. Eigene Darstellung.}
\label{fig:fftfluss}
\end{figure}

Bei jedem Sample wird eine vollständige $N$-Punkt-Transformation über die
letzten $N$ Samples berechnet. Es handelt sich also um eine gleitende
Fenster-FFT. Ein Block von $N$ Samples erfordert damit $N$ vollständige
Transformationen. Diese Eigenschaft ist bei der Interpretation der
Laufzeitwerte in Kapitel~6 zu berücksichtigen.

\subsubsection*{Wachstum des erzeugten Codes}

Faust erzeugt für diesen Baustein keine Schleife. Der Grund liegt im
Übersetzungsvorgang aus Abschnitt~3.5: Der Compiler wertet die
Blockdiagrammdefinitionen vollständig aus und überführt sie in eine explizite
Normalform. Jede einzelne Rechenoperation der Transformation steht danach als
eigener Ausdruck im erzeugten Code. Eine Radix-2-Transformation der Länge $N$
besteht aus $\tfrac{N}{2}\log_2 N$ Schmetterlingsoperationen. Der erzeugte
Code wächst daher etwa in dieser Größenordnung mit $N$.

Die Übersetzungszeit wächst noch stärker als der Code. Eine Verdopplung von
$N$ vervielfacht sie um ein Mehrfaches. Ab einer Länge von 512 überschreitet
die Übersetzung eine gesetzte Zeitgrenze von 30 Sekunden. Größere Längen
wurden deshalb nicht gemessen. Dieser Zusammenhang wird in Abschnitt~6.5
gesondert ausgewertet.

\subsubsection*{C++-Referenz}

Die C++-Referenz ist eine handgeschriebene iterative Radix-2-FFT nach
Cooley-Tukey mit Dezimation im Zeitbereich. Sie folgt dem in
Abschnitt~2.3.2 beschriebenen Zerlegungsprinzip. Die Umsetzung arbeitet in
drei Schritten. Zuerst werden die Eingangswerte nach ihrem bitweise
umgekehrten Index umsortiert. Anschließend folgen $\log_2 N$ Stufen. In jeder
Stufe werden Wertepaare mit einem festen Abstand miteinander verrechnet. Eine
solche Verrechnung heißt Schmetterlingsoperation und lautet

\begin{align}
a' &= a + w \cdot b, \\
b' &= a - w \cdot b,
\end{align}

wobei $w$ ein Drehfaktor ist. Der Abstand der Paare verdoppelt sich von Stufe
zu Stufe. Die Drehfaktoren und die Umsortierungstabelle werden bei der
Initialisierung einmalig berechnet und danach nur noch gelesen.

Die Referenz arbeitet ebenfalls als gleitende Fenster-FFT. Sie führt für
jedes Sample denselben vollständigen Ablauf aus wie die Faust-Variante.
Beide Implementierungen leisten damit dieselbe Arbeit.

Bewusst wurde keine etablierte FFT-Bibliothek wie FFTW eingesetzt. Solche
Bibliotheken enthalten handoptimierten SIMD-Code. Ein Vergleich gegen sie
würde nicht Faust gegen C++ messen, sondern Faust gegen eine spezialisierte
Bibliothek. Die Referenz verwendet außerdem dieselbe Fensterorientierung wie
Faust, bei der das neueste Sample an erster Stelle steht.


\section{Benchmark-Tools und Messgrößen}
\label{sec:messgroessen}

\subsection{Laufzeitmessung}

Alle Laufzeitmessungen erfolgen mit Google~Benchmark. Die Bibliothek wird als
einzige Laufzeit-Engine verwendet. Zwei Eigenschaften sprechen für diese
Wahl.

Erstens bestimmt die Bibliothek die Anzahl der Wiederholungen innerhalb einer
Messung selbst. Sie verlängert eine Messung so lange, bis die Gesamtdauer
deutlich über der Auflösung des Systemzeitgebers liegt. Einzelne sehr kurze
Aufrufe lassen sich dadurch überhaupt erst zuverlässig messen.

Zweitens weist sie die Prozessorzeit getrennt von der vergangenen Realzeit
aus. Die Realzeit wird mit einer gewöhnlichen Systemuhr bestimmt. Die
Prozessorzeit stammt dagegen von einer eigenen Uhr des Betriebssystems, die
nur läuft, während der Prozess tatsächlich auf einem Kern rechnet. Unter
Linux geschieht das über den Zeitgeber \texttt{CLOCK\_PROCESS\_CPUTIME\_ID}.
Wird der Prozess vom Betriebssystem verdrängt, steht diese Uhr still. Fremde
Systemlast verlängert daher die Realzeit, nicht aber die Prozessorzeit. Für
den Vergleich zweier Implementierungen ist deshalb nur die Prozessorzeit
aussagekräftig.

Jede Kombination aus Baustein, Implementierung und Blockgröße wird 50-mal
wiederholt gemessen. Die Einzelwerte der Wiederholungen bleiben in der
Ergebnisdatei erhalten. Alle berichteten Kennzahlen stammen dadurch aus
derselben Messung. Die zentrale Messgröße ist die Prozessorzeit pro Aufruf
der Verarbeitungsmethode, also pro Audioblock.
% OFFEN: Die Benchmarks melden ueber SetItemsProcessed zusaetzlich einen
% Durchsatz in Samples pro Sekunde. Dieser Wert steht in den JSON-Dateien,
% wird aber in keiner Abbildung und in keiner Summary-CSV ausgewertet.
% Entweder eine Auswertung ergaenzen oder den Durchsatz gar nicht erwaehnen.

\subsection{Berichtete Kennzahlen}

Als Hauptkennzahl dient der Median. Laufzeitverteilungen sind rechtsschief:
Einzelne Ausreißer durch Aktivität des Betriebssystems verlängern eine
Messung, verkürzen sie aber nie. Der Mittelwert wird dadurch verzerrt, der
Median nicht.

Ergänzend werden drei weitere Kennzahlen ausgewiesen. Der Mittelwert dient
der Vergleichbarkeit mit anderen Arbeiten. Der Interquartilsabstand
beschreibt die Streuung der mittleren Hälfte aller Messwerte. Das
95.~Perzentil beschreibt den nahezu ungünstigsten Fall: 95 Prozent aller
Blöcke wurden schneller verarbeitet als dieser Wert. Diese Kennzahl ist für
Audioanwendungen wichtiger als der Mittelwert. Verpasst ein einzelner Block
seine Frist, entsteht eine hörbare Störung. Ein guter Durchschnitt hilft in
diesem Fall nicht. Das Maximum wäre die naheliegende Alternative, wird aber
von einzelnen Störungen des Betriebssystems bestimmt und schwankt daher stark.
Das 95.~Perzentil ist gegenüber solchen Ausreißern unempfindlich.

\subsection{Darstellung der Messunsicherheit}

Die Streuung der Messungen wird auf drei Wegen sichtbar gemacht. In den
Laufzeitdiagrammen wird um den Median ein Band vom ersten bis zum dritten
Quartil eingezeichnet. In den Ergebnistabellen wird der Variationskoeffizient
angegeben, also die Standardabweichung geteilt durch den Mittelwert. Er macht
Messreihen unterschiedlicher Größenordnung vergleichbar. Für das Verhältnis
zwischen beiden Implementierungen wird zusätzlich ein Konfidenzintervall
mittels Bootstrap-Verfahren angegeben, da sich Unsicherheiten bei einem
Quotienten nicht einfach addieren lassen.
% OFFEN: Das Skript fuer Bootstrap-Intervall und Signifikanztest liegt im
% Repository, wird aber bisher nicht aufgerufen und erzeugt keine Datei.
% Diesen Absatz erst uebernehmen, wenn die Zahlen tatsaechlich vorliegen.

\subsection{Echtzeitbezug}

Die gemessenen Zeiten werden der verfügbaren Rechenzeit gegenübergestellt.
Bei nicht überlappender Blockverarbeitung entspricht diese nach
Gleichung~(2.6) der Blockdauer. Für 256 Samples bei 48\,kHz ergibt das
rechnerisch etwa 5{,}3\,ms.

Dieser Wert ist eine obere Schranke und nicht die in der Praxis verfügbare
Zeit. Dafür gibt es drei Gründe. Erstens verwenden Anwendungen mit geringer
Latenz meist kleinere Puffer von 64 bis 128 Samples, wodurch die Frist auf
etwa 1{,}3 bis 2{,}7\,ms sinkt. Zweitens gilt bei überlappender Verarbeitung
nach Gleichung~(2.7) die kürzere Hop-Zeit. Drittens teilt sich ein einzelner
DSP-Baustein den Verarbeitungsaufruf mit allen übrigen Bausteinen der
Signalkette. Ihm steht daher nur ein Bruchteil der Blockdauer zu. Der
Vergleich mit der Blockdauer zeigt somit, ob ein Baustein grundsätzlich
echtzeittauglich ist, nicht ob er in einer konkreten Anwendung ausreicht.

\subsection{Weitere Messgrößen}

Neben der Laufzeit werden vier weitere Größen erhoben.
Tabelle~\ref{tab:messgroessen} fasst sie zusammen.

\begin{table}[htb]
\centering
\caption{Erhobene Messgrößen und die verwendeten Werkzeuge}
\label{tab:messgroessen}
\begin{tabular}{lll}
\hline
Messgröße & Werkzeug & Einheit \\
\hline
Prozessorzeit pro Block & Google Benchmark & ns \\
Numerische Genauigkeit & Vergleich gegen NumPy-Referenz & absoluter Fehler \\
Speicherbedarf & Zustandsgröße und gezählte Allokationen & Byte \\
Binärgröße & Abschnittsgrößen der Objektdatei & Byte \\
Übersetzungszeit & wiederholter Aufruf des Faust-Compilers & s \\
\hline
\end{tabular}
\end{table}

\subsubsection*{Speicherbedarf}

Ein DSP-Baustein braucht Speicher an zwei Stellen. Ein Teil gehört fest zum
Objekt und steht schon zur Übersetzungszeit fest. Ein anderer Teil wird erst
zur Laufzeit angefordert. Faust legt seine Verzögerungsleitungen fest im
Objekt an und fordert zur Laufzeit nichts an. Die C++-Referenzen verwenden
dagegen dynamische Felder und fordern ihren Speicher zur Laufzeit an.

Ein Vergleich, der nur eine der beiden Stellen betrachtet, ist deshalb
wertlos. Gemessen werden daher beide: die feste Größe des Objekts und die
Menge des zur Laufzeit angeforderten Speichers. Berichtet wird ihre Summe.

Nicht gemessen wird der Speicherhöchststand des gesamten Testprogramms.
Dieser Wert wird von den Puffern des Testprogramms bestimmt und ist für alle
Bausteine nahezu gleich.

Zusätzlich wird gezählt, wie oft ein Baustein während der Verarbeitung
Speicher anfordert. Dieser Wert muss null sein. Eine Speicheranforderung
während der Verarbeitung kann beliebig lange dauern und macht einen Baustein
für Echtzeitanwendungen unbrauchbar.

\subsubsection*{Binärgröße}

Die Binärgröße wird nicht am fertigen Programm gemessen. Ein gelinktes
Programm enthält immer die Laufzeitbibliothek und den Startcode. Diese
konstante Grundlast ist um ein Vielfaches größer als der Baustein selbst und
überdeckt jeden Unterschied. Gemessen werden daher die Abschnittsgrößen einer
Übersetzungseinheit, die ausschließlich den jeweiligen Baustein enthält.
% OFFEN: Beide Messungen (Speicher, Binaergroesse) muessen mit dem neuen
% Verfahren noch einmal durchgefuehrt werden.


\section{Hardwareumgebung}
\label{sec:hardware}

Alle Messungen liefen auf demselben Rechner.
Tabelle~\ref{tab:umgebung} listet die Eigenschaften der Messumgebung auf. Die
Angaben werden vor jedem Messlauf automatisch protokolliert und zusammen mit
den Ergebnissen abgelegt.

\begin{table}[htb]
\centering
\caption{Mess- und Entwicklungsumgebung}
\label{tab:umgebung}
\begin{tabular}{ll}
\hline
Eigenschaft & Wert \\
\hline
Prozessor & AMD Ryzen 7 4700U \\
Kerne / Threads & 8 / 8 \\
Befehlssatzerweiterungen & SSE4.2, AVX, AVX2, FMA \\
L1-Datencache & 32 KiB je Kern \\
L2-Cache & 512 KiB je Kern \\
L3-Cache & 4 MiB gemeinsam \\
Betriebssystem & Ubuntu 24.04 unter WSL2 \\
Kernel & 6.18.33.2-microsoft-standard-WSL2 \\
Compiler & g++ 13.3.0 \\
Build-System & CMake 3.29.3, Ninja 1.11.1 \\
Faust-Compiler & 2.83.1 \\
Referenzumgebung & Python 3.12.3, NumPy 2.5.1 \\
\hline
\end{tabular}
\end{table}

\subsection{Bedeutung der Befehlssatzerweiterungen}

Zwei Einträge der Tabelle sind für die Auswertung wesentlich. AVX2 ist eine
Erweiterung des Befehlssatzes um Vektorbefehle. Ein solcher Befehl verarbeitet
mehrere Werte gleichzeitig. Die Register sind 256 Bit breit und fassen damit
acht Gleitkommazahlen einfacher Genauigkeit. Die automatische Vektorisierung
aus Abschnitt~2.4.1 kann dadurch im günstigsten Fall acht Samples pro Befehl
verarbeiten. FMA steht für die verschmolzene Multiplikation und Addition. Ein
FMA-Befehl berechnet $a \cdot b + c$ in einem Schritt mit nur einer Rundung.
Genau diese Rechenoperation bildet den Kern von FIR-Filtern und
Biquad-Filtern.

Der Prozessor unterstützt kein AVX-512. Die Vektorbreite von acht Werten ist
damit die obere Grenze für den Gewinn durch Vektorisierung. Diese Angabe ist
nötig, um die gemessenen Unterschiede einordnen zu können. Ohne sie ließe
sich nicht beurteilen, ob ein beobachteter Faktor nahe am technisch Möglichen
liegt.

\subsection{Bedeutung der Cachegrößen}

Die Laufzeit eines DSP-Bausteins hängt davon ab, ob seine Daten im schnellen
Cache liegen. Muss ein Filter seine Koeffizienten aus dem Hauptspeicher
holen, bestimmt die Speicherzugriffszeit das Ergebnis. Der Unterschied
zwischen zwei erzeugten Codevarianten wäre dann nicht mehr erkennbar. Für die
Aussagekraft der Messung ist deshalb zu prüfen, ob die Daten in den
L1-Datencache passen.

Der längste untersuchte FIR-Filter benötigt 256 Koeffizienten und eine
Verzögerungsleitung derselben Länge. Bei einfacher Genauigkeit mit 4 Byte je
Wert ergibt das
$256 \cdot 4\,\mathrm{B} + 256 \cdot 4\,\mathrm{B} = 2048\,\mathrm{B}$, also
2 KiB. Das sind etwa 6 Prozent des L1-Datencaches von 32 KiB. Alle
FIR-Varianten arbeiten damit vollständig im schnellsten Cache. Die Messung
erfasst daher den erzeugten Rechencode und nicht die Speicheranbindung.

\subsection{Einschränkungen der Umgebung}

Die Messungen liefen unter WSL2 und damit in einer virtualisierten Umgebung.
Diese Umgebung erlaubt keinen Zugriff auf die Frequenzsteuerung des
Prozessors. Die Taktfrequenz ist während der Messungen deshalb nicht
festgelegt. Absolute Zeitangaben sind aus diesem Grund nur eingeschränkt
übertragbar. Der Vergleich zweier Implementierungen bleibt davon jedoch
weitgehend unberührt, da beide Implementierungen innerhalb desselben
Programmlaufs und unmittelbar nacheinander gemessen werden. Vor jedem
Messlauf wird zusätzlich die Systemlast protokolliert. Die Auswirkungen
dieser Einschränkungen werden in Abschnitt~7.2 diskutiert.


\section{Compilerkonfigurationen}
\label{sec:compilerkonfig}

\subsection{Grundsätze}

Für den Vergleich gelten drei Regeln. Erstens werden beide Implementierungen
immer mit denselben Compilerflags übersetzt. Zweitens liegen beide
Implementierungen im selben Programm und werden im selben Lauf gemessen.
Drittens verwendet die C++-Referenz keine expliziten SIMD-Befehle.

Der Faust-Compiler wird ohne Architekturdatei aufgerufen. Übergeben wird nur
der Name der zu erzeugenden Klasse. Faust arbeitet damit im skalaren Modus
nach Abschnitt~3.6.1. Der Vektor-Modus und der Parallel-Modus sind nicht
Gegenstand der Messungen. Ebenso wird durchgehend einfache Genauigkeit
verwendet.

\subsection{Aktivierung der automatischen Vektorisierung}

Die automatische Vektorisierung ist kein eigener Schalter, sondern an die
Optimierungsstufe gekoppelt. Bei GCC aktiviert \texttt{-O3} die
Vektorisierung von Schleifen. Ab Version~12 ist sie bereits bei \texttt{-O2}
aktiv, dort allerdings mit einem zurückhaltenden Kostenmodell. Die Option
\texttt{-march=native} erlaubt dem Compiler zusätzlich, die AVX2- und
FMA-Befehle des Testrechners tatsächlich zu verwenden. Ohne diese Option
beschränkt er sich auf einen älteren Basisbefehlssatz. Beide Angaben
zusammen sind notwendig: \texttt{-O3} schaltet die Vektorisierung ein,
\texttt{-march=native} stellt ihr die breiten Befehle zur Verfügung. Weitere
Einstellungen wurden nicht vorgenommen, insbesondere keine Pragmas und keine
Intrinsics.
% OFFEN: Die Aussage zu den Standardeinstellungen von GCC auf dem
% Messrechner belegen mit:  g++ -O2 -Q --help=optimizers | grep vectorize
% und  g++ -O3 -Q --help=optimizers | grep vectorize
% Ergebnis in einer Fussnote oder im Anhang festhalten.

\subsection{Flag-Matrix}

Als Referenzkonfiguration dient \texttt{-O3 -march=native -ffast-math
-DNDEBUG}. Zusätzlich wird eine Matrix aus sieben Konfigurationen gemessen.
Tabelle~\ref{tab:flags} zeigt sie zusammen mit ihrer jeweiligen Begründung.

\begin{table}[htb]
\centering
\caption{Untersuchte Compilerkonfigurationen}
\label{tab:flags}
\begin{tabular}{lll}
\hline
Konfiguration & Zweck & Verwendet für \\
\hline
\texttt{-O0} & unoptimierter Ausgangszustand & Laufzeit \\
\texttt{-O1} & einfache lokale Optimierungen & Laufzeit \\
\texttt{-O2} & übliche Stufe für Produktivcode & Laufzeit \\
\texttt{-O3} & Entrollung und Vektorisierung & Laufzeit \\
\texttt{-O3 -march=native} & Nutzung von AVX2 und FMA & Laufzeit \\
\texttt{-O3 -march=native -ffast-math} & Umordnung von Gleitkommaoperationen & Laufzeit \\
\texttt{-Os} & Optimierung auf Codegröße & Binärgröße \\
\hline
\end{tabular}
\end{table}

Die letzte Zeile betrifft ausschließlich die Messung der Binärgröße. Die
Referenzkonfiguration optimiert bewusst auf Geschwindigkeit. Sie entrollt
Schleifen und erzeugt Vektorbefehle, was den Code vergrößert. Eine
Binärgröße, die unter diesen Bedingungen gemessen wird, beantwortet die
Teilfrage aus Abschnitt~1.2 nur unvollständig. Binärgröße ist vor allem für
eingebettete Systeme von Bedeutung, und dort wird auf Codegröße optimiert.

Die Option \texttt{-Os} schaltet Entrollung und Vektorisierung ab. Sie wirkt
dabei ausschließlich auf den C++-Compiler. Was der Faust-Compiler bereits
ausgerollt hat, kann der C++-Compiler nicht wieder in eine Schleife
zurückführen. Der zusätzliche Messlauf trennt damit zwei Ursachen
voneinander: Unter der Referenzkonfiguration mischen sich die Struktur des
erzeugten Codes und die Entrollung durch den C++-Compiler. Unter \texttt{-Os}
bleibt nur die Struktur übrig. Der Unterschied, der dort noch sichtbar ist,
geht auf die Codegenerierung von Faust zurück.

Für die Laufzeitmessungen wird \texttt{-Os} nicht verwendet, da eine auf
Größe optimierte Übersetzung den Laufzeitvergleich verfälschen würde.

Jede Konfiguration enthält zusätzlich das Makro \texttt{NDEBUG}. Das ist
notwendig, weil Google~Benchmark den Übersetzungsmodus nicht am Build-Typ des
Build-Systems erkennt, sondern an diesem Makro. Fehlt es, meldet die
Bibliothek einen Debug-Build, und die gemessenen Zeiten sind nicht
verwertbar.

Die Stufen \texttt{-O0} und \texttt{-O1} sind nicht als praxisnahe
Konfigurationen zu verstehen. Sie dienen dazu, den Einfluss der Optimierung
auf den Vergleich sichtbar zu machen. Die Option \texttt{-ffast-math} erlaubt
dem Compiler, Gleitkommaoperationen umzuordnen. Sie kann dadurch die Laufzeit
verändern und gleichzeitig das Rechenergebnis beeinflussen. Beide Wirkungen
sind Gegenstand der Auswertung in Abschnitt~6.2.

Die tatsächlich verwendeten Flags werden bei jedem Build maschinenlesbar
abgelegt und vor der Messung automatisch überprüft. Damit ist nachweisbar,
dass beide Implementierungen unter identischen Bedingungen übersetzt wurden.

Die Programme zur Genauigkeitsprüfung bilden eine Ausnahme. Sie werden
zusätzlich mit \texttt{-fno-fast-math} übersetzt. Sie sollen die tatsächliche
Rundung in einfacher Genauigkeit messen und nicht das Ergebnis einer
Umordnung durch den Compiler.
% OFFEN: Damit werden Korrektheit und Laufzeit an unterschiedlich
% uebersetzten Programmen gemessen. Entweder zusaetzlich einen
% Genauigkeitslauf mit -ffast-math erheben oder den Punkt in Abschnitt 7.3
% als Einschraenkung der Konstruktvaliditaet aufnehmen.


\section{Testdatengenerierung}
\label{sec:testdaten}

Alle Testsignale werden mit einem Python-Skript erzeugt und als Rohdaten in
einfacher Genauigkeit abgelegt. Der Zufallszahlengenerator verwendet einen
festen Startwert. Die Signale sind dadurch über beliebig viele Läufe hinweg
identisch. Es werden zwei Gruppen von Signalen unterschieden.

\subsection{Signale für die Laufzeitmessung}

Für die Laufzeitmessungen dient gleichverteiltes Rauschen im Wertebereich von
$-1$ bis $1$. Für jede Blockgröße liegt eine eigene Datei vor. Die Benchmarks
laden diese Datei zu Beginn der Messung und verarbeiten sie anschließend
wiederholt.

Das Laden aus einer Datei ist eine bewusste Entscheidung. Würde das
Eingangssignal im Programm selbst als Konstante erzeugt, könnte der Compiler
Teile der Berechnung bereits zur Übersetzungszeit auflösen. Die Messung würde
dann nicht mehr die Verarbeitung erfassen, sondern nur noch das Ergebnis
dieser Optimierung. Zusätzlich verhindern zwei Anweisungen der
Benchmark-Bibliothek, dass der Compiler die Berechnung als unbenutzt
einstuft und entfernt.

\subsection{Signale für die Genauigkeitsprüfung}

Für die Genauigkeitsprüfung werden sieben Signaltypen in vier Längen erzeugt,
nämlich 64, 256, 1024 und 4096 Samples. Tabelle~\ref{tab:signale} beschreibt
die Typen und ihren jeweiligen Zweck. Die Auswahl deckt sowohl typische
Audiosignale als auch Randfälle ab.

\begin{table}[htb]
\centering
\caption{Signaltypen für die Genauigkeitsprüfung}
\label{tab:signale}
\begin{tabular}{lll}
\hline
Typ & Definition & Zweck \\
\hline
Impuls & $x[0]=1$, sonst 0 & Impulsantwort \\
Sprung & konstant 1 & Sprungantwort, Einschwingen \\
Stille & konstant 0 & trivialer Randfall \\
Gleichanteil & konstant 0{,}5 & Linearität \\
Sinus & 1 kHz bei 48 kHz, Amplitude 0{,}8 & typisches Audiosignal \\
Extremwerte & alternierend $\pm 1$ & ungünstigster Fall für Rundung \\
Rauschen & gleichverteilt, eigener Startwert & Durchschnittsverhalten \\
\hline
\end{tabular}
\end{table}

\subsection{Referenzwerte}

Ein Vergleich zwischen Faust und der C++-Referenz allein wäre nicht
ausreichend. Machen beide Implementierungen denselben systematischen Fehler,
bliebe dieser unentdeckt. Daher wird zusätzlich eine Referenz in Python
berechnet. Sie verwendet dieselben Algorithmen, rechnet jedoch in doppelter
Genauigkeit. Beide Implementierungen werden gegen diese Referenz geprüft. Die
Referenzwerte werden ebenfalls als Rohdaten abgelegt und vor jeder
Genauigkeitsprüfung neu erzeugt.


% ===========================================================================
%  OFFENE PUNKTE - vor der Abgabe abarbeiten und diese Liste loeschen
%
%  1. Abschnitt 4.1: Akkumulator auch in 1.2 und 2.2 einfuehren.
%  2. Abschnitt 4.1.2: Regler gegen Konstante bei der Verstaerkungsstufe.
%  3. Abschnitt 4.2: Durchsatz wird gemessen, aber nirgends ausgewertet -
%     entweder Auswertung ergaenzen oder Erwaehnung streichen.
%  4. Abschnitt 4.2: Bootstrap-Intervall erst erwaehnen, wenn erhoben.
%  5. Abschnitt 4.2: Speicher- und Binaergroessenmessung neu erheben.
%  6. Abschnitt 4.4: GCC-Vektorisierungsdefaults auf dem Messrechner belegen.
%  7. Abschnitt 4.4: Genauigkeit zusaetzlich mit -ffast-math oder als
%     Einschraenkung in 7.3 aufnehmen.
%  8. Zitierschluessel an die eigene .bib anpassen. Fuer Google Benchmark
%     existiert bisher kein Eintrag im Quellenverzeichnis.
%  9. Schreibweise vereinheitlichen: Kapitel 2 verwendet "Faust",
%     Kapitel 3 verwendet "FAUST". Dieses Kapitel verwendet "Faust".
% ===========================================================================